% Problema 2, parte 2: campi da incollare nella piattaforma (uno per sezione)

% ===== CAMPO: 1. Risultato e ambito =====
\(U(Q_5) = 88 = |E(Q_5)| + 8\). Etichettatura ottima (ordine crescente
di etichetta; posizione \(i\) della stringa = coordinata \(i\);
convenzione irrilevante perché le permutazioni di coordinate sono
automorfismi di \(Q_5\)), con 8 valli:
\texttt{00000,\ 11110,\ 11001,\ 10101,\ 01101,\ 10011,\ 01011,\ 00111}
sono le etichette 1--8 (le valli), seguite dalle foglie e dai picchi
come nella lista completa in §2 (tabella dell'upper bound). La lista
completa a 32 elementi è prodotta e verificata da
\texttt{certificati/costruzione\_88.py} e
\texttt{certificati/verifica\_dfs\_88.py}.


% ===== CAMPO: 2. Dimostrazione =====
\emph{(Testo del tentativo
\texttt{runs/p2\_q5/attempts/attempt\_001.json} del Researcher, riletto
lemma per lemma dal team; le verifiche computazionali sono state
rieseguite dall'orchestratore in
\texttt{runs/p2\_q5/verifica/attempt\_001/}.)}

\paragraph{\texorpdfstring{\(U(Q_5) = 88\)}{U(Q\_5) = 88}}\label{uq_5-88}

\textbf{Conventions.} \(V(Q_5)=\{0,1\}^5\), \(u\sim v\) iff they differ
in exactly one coordinate; \(|V|=32\), \(|E|=80\), every vertex has
degree \(5\). In a 0/1 string the \(i\)-th character from the left is
coordinate \(i\) (irrelevant for the count: coordinate permutations are
automorphisms). For a labelling \(f\), orient each edge from the smaller
to the larger label (\(u\to v\) iff \(uv\in E\), \(f(u)<f(v)\));
\(\deg^-(v)=\#\{u:u\to v\}\), \(\deg^+(v)=\#\{w:v\to w\}\),
\(\deg^-(v)+\deg^+(v)=5\). Let \(p(v)\) be the number of uphill paths
whose last vertex is \(v\). A \emph{valley} is a vertex with
\(\deg^-=0\); a \emph{peak} is a vertex with \(\deg^+=0\). The vertex
with label \(1\) is a valley, so \(v:=\#\text{valleys}\ge1\).

\subparagraph{0. Preliminaries (restated with proof for
self-containedness)}\label{preliminaries-restated-with-proof-for-self-containedness}

\textbf{Lemma 1 (recursion).}
\(p(v)=[v\text{ valley}]+\sum_{u\to v}p(u)\) and the total number of
uphill paths is \(T=\sum_v p(v)\).

\emph{Proof.} An uphill path ending at \(v\) has length \(1\) --- then
it is \((v)\) and \(v\) is a valley; conversely a valley gives exactly
this path --- or length \(\ge2\); deleting \(v\) gives an uphill path
ending at its penultimate vertex \(u\) with \(u\to v\); conversely
appending \(v\) to an uphill path ending at \(u\) with \(u\to v\) gives
an uphill path ending at \(v\) (the increasing condition is exactly
\(f(u)<f(v)\)). These maps are mutually inverse and different \(u\) give
disjoint sets. Every uphill path has a unique last vertex. \(\square\)

\textbf{Lemma 2.} For every vertex \(v\): \(p(v)\ge1\), and
\(p(v)\ge\deg^-(v)\).

\emph{Proof.} From \(v\) move repeatedly to a neighbour with smaller
label while one exists; labels strictly decrease, so this stops at a
valley; the reversed walk is an uphill path ending at \(v\). Hence
\(p(v)\ge1\) for all \(v\), and by Lemma 1,
\(p(v)\ge\sum_{u\to v}p(u)\ge\deg^-(v)\). \(\square\)

\subparagraph{1. The excess identity}\label{the-excess-identity}

Define the \textbf{excess} \(X:=\sum_{(u\to w)\in E}\bigl(p(u)-1\bigr)\)
(sum over all \(80\) oriented edges). By Lemma 2 every term is \(\ge0\).

\textbf{Lemma 3.} \(T=|E|+v+X = 80+v+X\).

\emph{Proof.} By Lemma 1,
\(T=\sum_v p(v)=v+\sum_{w}\sum_{u\to w}p(u)=v+\sum_{(u\to w)\in E}\bigl(1+(p(u)-1)\bigr)=v+|E|+X\),
since each edge is counted once, at its head. \(\square\)

Grouping the terms of \(X\) by tail: \(X=\sum_u \deg^+(u)\,(p(u)-1)\).
Call \(u\) \textbf{heavy} if \(\deg^+(u)\ge1\) and \(p(u)\ge2\); let
\(H\) be the set of heavy vertices and \(c(u):=\deg^+(u)(p(u)-1)\) its
contribution, so \(X=\sum_{u\in H}c(u)\). Vertices with \(p=1\) are
called \textbf{light}.

\textbf{Lemma 4 (structure of light vertices).} If \(p(u)=1\) then
either \(u\) is a valley or \(\deg^-(u)=1\) and its unique in-neighbour
is light. Consequently, for any set \(L\) of light vertices, the
subgraph of \(Q_5\) induced by \(L\) is acyclic.

\emph{Proof.} By Lemma 2, \(\deg^-(u)\le p(u)=1\). If \(\deg^-(u)=1\)
with in-neighbour \(u'\), Lemma 1 gives \(1=p(u)=p(u')\). For the second
claim: a cycle inside \(L\) contains a vertex of maximal label, which
has two in-neighbours on the cycle, contradicting \(\deg^-\le1\).
\(\square\)

\textbf{Lemma 5 (every heavy vertex costs at least 3).} For every heavy
\(u\), \(c(u)\ge3\). Moreover, writing \(i=\deg^-(u)\) (so
\(\deg^+(u)=5-i\), \(1\le i\le4\) since \(\deg^+\ge1\) and \(p\ge2\)
excludes valleys): - if \(i=1\): \(p(u)=p(u')\ge2\) for the in-neighbour
\(u'\), which is then itself heavy (\(\deg^+(u')\ge1\) because
\(u'\to u\)), and \(c(u)=4(p(u)-1)\ge4\); - if \(i=2\):
\(c(u)=3(p(u)-1)\ge3\), with equality iff \(p(u)=2\) iff both
in-neighbours are light; - if \(i=3\): \(c(u)=2(p(u)-1)\ge4\) (as
\(p(u)\ge3\)), with equality iff all three in-neighbours are light; - if
\(i=4\): \(c(u)=p(u)-1\ge3\), with equality iff all four in-neighbours
are light.

\emph{Proof.} Direct from Lemma 2 (\(p(u)\ge i\)) and Lemma 1
(\(p(u)=\sum_{u'\to u}p(u')\) with each \(p(u')\ge1\); equality
\(p(u)=i\) iff every in-neighbour has \(p=1\)). Note also: if an
in-neighbour \(u'\) of a heavy \(u\) has \(p(u')\ge2\) then \(u'\) is
heavy (it has \(\deg^+\ge1\)). \(\square\)

\textbf{Corollary 6.} \(X\in\{0\}\cup[3,\infty)\), and \(X\le 7\) forces
\(|H|\le2\).

\textbf{Lemma 7 (in-degree count).} Let \(P\) be the set of peaks,
\(q=|P|\), \(s_H=\sum_{u\in H}\deg^-(u)\). If every non-heavy non-peak
vertex is light (which holds automatically for every vertex with
\(\deg^+\ge1\) that is not heavy, by definition of heavy), then
\[4q=48+v+|H|-s_H .\]

\emph{Proof.} Peaks have \(\deg^-=5\); no peak is heavy (heavy needs
\(\deg^+\ge1\)) and no valley is heavy. The vertices outside \(P\cup H\)
are light: valleys (\(\deg^-=0\), \(v\) of them) and light non-valleys
(\(\deg^-=1\) by Lemma 4, \(32-q-|H|-v\) of them). Summing in-degrees,
\(80=\sum_u\deg^-(u)=5q+s_H+(32-q-|H|-v)\), i.e.~\(4q=48+v+|H|-s_H\).
\(\square\)

Also note: \textbf{peaks form an independent set} (if two peaks were
adjacent, the one with the smaller label would have \(\deg^+\ge1\)).

\subparagraph{\texorpdfstring{2. Case analysis: every labelling with
\(T\le87\) has one of four
structures}{2. Case analysis: every labelling with T\textbackslash le87 has one of four structures}}\label{case-analysis-every-labelling-with-tle87-has-one-of-four-structures}

Assume \(T\le87\), i.e.~(Lemma 3) \(v+X\le7\) with \(v\ge1\), so
\(X\le6\). By Corollary 6, \(X\in\{0,3,4,5,6\}\) and \(|H|\le2\).

\textbf{Case \(|H|=2\).} Then \(X\ge6\), so \(X=6\), \(v=1\), and both
heavy vertices have \(c=3\); by Lemma 5 each has
\((\deg^-,p)\in\{(2,2),(4,4)\}\), so \(s_H\in\{4,6,8\}\) and Lemma 7
gives \(4q=48+1+2-s_H\in\{47,45,43\}\), none divisible by \(4\).
\textbf{Impossible.}

\textbf{Case \(|H|=0\).} Then \(X=0\), \(T=80+v\), and Lemma 7 gives
\(4q=48+v\), so \(v\equiv0\pmod 4\); with \(T\le87\) this forces
\(v=4\), \(q=13\), \(T=84\). Every non-peak vertex is light; so
\(F:=V\setminus P\) (19 vertices) induces an acyclic subgraph (Lemma 4),
and \(P\) is an independent set of size \(13\). \textbf{(Structure A.)}

\textbf{Case \(|H|=1\), \(H=\{u\}\).} Then \(X=c(u)\le6\). By Lemma 5,
if some in-neighbour of \(u\) had \(p\ge2\) it would be heavy,
contradicting \(|H|=1\); hence all in-neighbours of \(u\) are light and
\(p(u)=\deg^-(u)=:i\); \(i=1\) is excluded (it needs a heavy
in-neighbour). So \((i,c(u))\in\{(2,3),(3,4),(4,3)\}\), and Lemma 7 with
\(|H|=1\), \(s_H=i\) gives \(4q=49+v-i\): - \(i=2\): \(4q=47+v\),
\(v\equiv1\pmod4\), and \(v+X=v+3\le7\) gives \(v=1\), \(q=12\),
\(T=84\). \textbf{(Structure B.)} - \(i=3\): \(4q=46+v\),
\(v\equiv2\pmod4\), \(v+4\le7\) gives \(v=2\), \(q=12\), \(T=86\).
\textbf{(Structure C.)} - \(i=4\): \(4q=45+v\), \(v\equiv3\pmod4\),
\(v+3\le7\) gives \(v=3\), \(q=12\), \(T=86\). \textbf{(Structure D.)}

In each of B, C, D: the out-neighbours \(w\) of \(u\) satisfy
\(p(w)\ge p(u)\ge2\) (Lemma 1), so if such a \(w\) had \(\deg^+(w)\ge1\)
it would be heavy; hence \textbf{all out-neighbours of \(u\) are peaks}.
The in-neighbours of \(u\) are non-peaks. All vertices other than \(u\)
and the peaks are light, so by Lemma 4 the set
\(L:=V\setminus(P\cup\{u\})\) induces an acyclic subgraph; \(P\) is an
independent set of size \(12\) and \(u\notin P\).

\textbf{Summary.} If \(T\le 87\) then there exist an independent set
\(I\subseteq V(Q_5)\) and a vertex \(u\notin I\) such that
\(V\setminus(I\cup\{u\})\) induces an acyclic subgraph, where either
\(|I|=12\) (structures B, C, D with \(I=P\)), or \(|I|=13\) and
\(V\setminus I\) is acyclic (structure A) --- in which case, choosing
any \(u_0\in I\) and \(I'=I\setminus\{u_0\}\), the pair \((I',u_0)\) has
\(|I'|=12\), \(I'\) independent, \(u_0\notin I'\), and
\(V\setminus(I'\cup\{u_0\})=V\setminus I\) acyclic. So in all cases:

\begin{quote}
\textbf{(★)} there is an independent set \(I\) of \(Q_5\) with
\(|I|=12\) and a vertex \(u\notin I\) such that \(Q_5-(I\cup\{u\})\) is
a forest.
\end{quote}

\subparagraph{\texorpdfstring{3. Exact finite verification: (★) never
holds in
\(Q_5\)}{3. Exact finite verification: (★) never holds in Q\_5}}\label{exact-finite-verification-never-holds-in-q_5}

\textbf{Reduction.} The maps \(x\mapsto x\oplus t\) are automorphisms of
\(Q_5\) and act transitively on vertices, and (★) is invariant under
automorphisms. Hence it suffices to check all independent \(12\)-sets
\(I\) containing the vertex \(00000\) and all \(u\notin I\).

\textbf{Computation 1} (\texttt{verifica\_q5\_indipendenti.py}, exact
integer arithmetic): backtracking enumerates all independent sets of
size \(12\) containing \(0\) in increasing vertex order (a vertex may be
added iff it is neither in the set nor adjacent to it), and for each of
the \(20\) vertices \(u\notin I\) tests acyclicity of the induced
subgraph on the \(19\) remaining vertices by the exact criterion
\emph{\#edges = \#vertices − \#components}. Output:
\texttt{insiemi\ indipendenti\ di\ taglia\ 12\ contenenti\ 0:\ 1377;\ coppie\ (I,u)\ esaminate:\ 27540;\ complementi\ aciclici\ trovati:\ 0;\ tempo\ 0.24s}.

\textbf{Computation 2} (\texttt{verifica\_q5\_bipartita.py}, independent
method, no symmetry reduction): every independent set is \(A\cup B\)
with \(A\) a subset of the 16 even-weight vertices and \(B\) a subset of
the odd-weight vertices not adjacent to \(A\); the script scans all
\(2^{16}\) bitmasks \(A\), all \(B\) of size \(12-|A|\), and tests
acyclicity by union--find over the edge list. Output:
\texttt{insiemi\ indipendenti\ di\ taglia\ 12\ (tutti):\ 3672;\ coppie\ (I,u)\ esaminate:\ 73440;\ complementi\ aciclici:\ 0;\ tempo\ 0.39s}.
Consistency: \(3672\cdot 12/32=1377\), as vertex-transitivity predicts.

Hence (★) is false, so no labelling of \(Q_5\) has \(T\le87\):
\[U(Q_5)\ge 88.\]

(The counts also make the impossibility of the \(X=0\), \(v=4\)
structure transparent by hand: an independent \(12\)- or \(13\)-set
turns out to lie in one parity class except for a few mixed sets, and
any three even vertices contain two at distance \(2\), which share two
odd neighbours and thus close a \(4\)-cycle in the complement; but the
proof relies only on the exhaustive check above.)

\subparagraph{4. Upper bound: a labelling with exactly 88 uphill
paths}\label{upper-bound-a-labelling-with-exactly-88-uphill-paths}

Let \(R=\{00000,11110\}\) and let \(P\) be the \(14\) even-weight
vertices other than \(R\). \(P\) is independent (even-weight vertices
are pairwise non-adjacent). The complement
\(F=R\cup\{\text{16 odd vertices}\}\) induces exactly the \(10\) edges
from \(R\) to the odd vertices (the two vertices of \(R\) are at
distance \(4\), so their neighbourhoods
\(\{10000,01000,00100,00010,00001\}\) and
\(\{01110,10110,11010,11100,11111\}\) are disjoint): \(F\) is a forest
with \(18\) vertices, \(10\) edges and \(8\) components, namely the two
stars centred at \(R\) and the \(6\) isolated odd vertices
\(11001,10101,01101,10011,01011,00111\).

Labelling (increasing label order; coordinate 1 on the left):

\begin{verbatim}
00000,11110,11001,10101,01101,10011,01011,00111,
10000,01000,00100,11100,00010,11010,10110,01110,00001,11111,
11000,10100,01100,10010,01010,00110,10001,01001,00101,11101,00011,11011,10111,01111
\end{verbatim}

i.e.~labels \(1\)--\(2\): the star centres; \(3\)--\(8\): the six
isolated odd vertices; \(9\)--\(18\): the ten star leaves;
\(19\)--\(32\): the fourteen peaks \(P\).

\emph{Count.} The \(8\) vertices with labels \(1\)--\(8\) have all
neighbours higher (their \(F\)-neighbours, if any, are leaves labelled
\(9\)--\(18\); their other neighbours are in \(P\)): they are valleys,
\(p=1\). Each leaf (labels \(9\)--\(18\)) has exactly one lower
neighbour, its star centre, so \(p=1\). Each peak has all \(5\)
neighbours lower, all in \(F\) with \(p=1\), so \(p=5\). Total
\(T=18\cdot1+14\cdot5=88\). Verified exactly by two independent scripts:
\texttt{costruzione\_88.py} (recursion of Lemma 1) and
\texttt{verifica\_dfs\_88.py} (explicit DFS enumeration of all uphill
paths from the \(8\) valleys), both printing \(88\) and the same \(8\)
valleys.

\subparagraph{5. Conclusion}\label{conclusion}

\(U(Q_5)=88=|E|+8\), attained by the labelling in §4. Consistency checks
(not part of the proof): the same Lemmas 3--7 for \(Q_4\) give
\(c(u)\ge2\) and \(3q=16+v\), reproducing \(U(Q_4)=34\) with \(v=2\);
for \(Q_3\) they reproduce \(U(Q_3)=14\). The identity \(T=80+v+X\) and
the fact \(X\in\{0\}\cup[3,\infty)\) were also sanity-checked on
\(200{,}000\) random labellings of \(Q_5\) (\texttt{sanity\_eccesso.py};
the smallest positive \(X\) seen was \(14\)). Simulated annealing over
labellings (\texttt{ricottura\_q5.py}, 200k swap moves per seed, exact
scoring) never went below \(88\) in any seed run so far (evidence only).


% ===== CAMPO: 3. Verifica: istruzioni, dipendenze, tempi =====
Dipendenze: Python 3 standard, nessun pacchetto. Tutti gli script usano
aritmetica intera esatta.

\begin{verbatim}
cd problema-2/certificati
python3 costruzione_88.py                 # costruisce l'etichettatura e conta 88 cammini (< 0.1 s)
python3 verifica_dfs_88.py                # conteggio indipendente per DFS esplicita (< 0.1 s)
python3 verifica_q5_indipendenti.py       # (★) confutata: 1377 insiemi con 0, 27540 coppie, 0 aciclici (0.47 s)
python3 verifica_q5_bipartita.py          # (★) confutata senza riduzione per simmetria: 3672 / 73440 / 0 (0.71 s)
python3 q5_stella_verifica_indipendente.py  # (★) confutata con metodo diverso (potatura delle foglie): 3672 / 73440 / 0 (0.40 s)
\end{verbatim}

Tempi misurati su portatile (riesecuzione fidata dell'orchestratore,
2026-09-26). Codice del tentativo: - \texttt{code\_1} (python, rigore
\texttt{exact}): Exact count of uphill paths of a labelling of Q\_d via
the recursion p(v) = {[}valley{]} + sum of p over lower neighbours
(shared module). - \texttt{code\_2} (python, rigore \texttt{exact}):
Certificate 1 for the lower bound: enumerate all independent 12-sets of
Q\_5 containing vertex 0 (1377 sets, symmetry reduction by XOR
translations) and, for eac - \texttt{code\_3} (python, rigore
\texttt{exact}): Certificate 2 for the lower bound, independent method
and no symmetry reduction: all independent 12-sets as A ∪ B (A ⊆ even
class over all 2\^{}16 bitmasks, B ⊆ od - \texttt{code\_4} (python,
rigore \texttt{exact}): Builds the 88-path labelling (peaks = even
vertices minus 00000, 11110) and counts exactly with the recursion:
prints 88, the 8 valleys and the 32-vertex order. - \texttt{code\_5}
(python, rigore \texttt{exact}): Independent verification of the
submitted labelling by explicit DFS enumeration of every uphill path
from each valley (string-based, no shared code): prints 88. -
\texttt{code\_6} (python, rigore \texttt{exact}): Sanity check (evidence
only, not part of the proof) of the identity T = 80 + \#valleys + X and
of X ∈ \{0\} ∪ {[}3,∞) on 200000 random labellings of Q\_5 (all asserts
- \texttt{code\_7} (python, rigore \texttt{float\_exploration\_only}):
Simulated annealing over labellings (exploration only): best value found
88 in every seed run (seeds 0-5 in a first run, seeds 0-3 in a second
run; 200k swap mo


% ===== CAMPO: 4. Fonti e contributo =====
\begin{itemize}
\tightlist
\item
  Lemma di conteggio \(p(v)=[\text{valle}]+\sum_{u\to v}p(u)\) e bound
  \(\#\ge|E|+\#\text{valli}\): nostri (parte 1, ridimostrati qui).
\item
  Contesto: IMO 2022 P6 (griglia), citato solo come contesto; nessun
  risultato preso dalla letteratura.
\item
  Ricerca arXiv deterministica (\texttt{runs/p2\_q5/literature.json},
  \texttt{problema-2/fonti/arxiv.json}): nessun lavoro sui cammini in
  salita sull'ipercubo; il problema risulta inedito. Contributo
  interamente nostro.
\end{itemize}


% ===== CAMPO: 5. Limiti e parti irrisolte =====
\begin{itemize}
\tightlist
\item
  Gap dichiarati dal Researcher nel tentativo:
\item
  The lower bound is computer-assisted at one point: statement (★) (no
  independent 12-set I and vertex u with Q\_5 − (I ∪ \{u\}) acyclic) is
  proved by exhaustive enumeration, not by hand. The reduction to (★) is
  fully by hand; the enumeration is exact, covers the whole finite set
  (all 3672 independent 12-sets × 20 outside vertices), runs in
  \textless{} 1 s and was done by two methods. A hand proof of (★) was
  only sketched (parity-class argument), not written out.
\item
  In Lemma 7 / §2 I use that every non-peak, non-heavy vertex is light;
  this is immediate from the definitions (a vertex with deg\^{}+ ≥ 1 and
  p ≥ 2 is heavy by definition) but the Referee should check that no
  case escapes: valleys (p=1), peaks (deg\^{}+=0), heavy, light ---
  these four classes cover V.
\item
  The step `out-neighbours of the unique heavy vertex u are peaks'
  (needed for nothing in the final finite check, since (★) only uses I
  independent of size 12 and u ∉ I) is included for completeness; the
  final reduction deliberately uses only the weaker necessary
  conditions, so any slip there does not affect the proof.
\item
  Simulated annealing results are exploration only and play no role in
  the proof; the second SA run (20 seeds) was still running when this
  report was written (4 seeds completed, all at 88).
\item
  Il giudice B ha classificato tutte le osservazioni su fonti e calcoli
  come INFO (fonti di contesto non caricate e non necessarie; versione
  di Python non registrata: 3.14, libreria standard; tempi riprodotti
  0.44--0.71 s). Nei primi due run l'unico punto MISSING riguardava le
  regole ufficiali, poi passate alla lettera dal preambolo
  dell'enunciato.
\item
  La cella è ``checked instantly'' sul valore: il rischio residuo è un
  errore nella riduzione a (★); per questo la riduzione è stata riletta
  a mano e (★) confutata da tre programmi diversi.
\end{itemize}

